%% ma: L10-B13-0008
%% lop: 10
%% chuong: 5
%% bai: 13
%% dang: 10.13.3
%% loai: TL
%% muc: NB
%% dap_an: "-"
%% nguon: "chuyen-de-cac-so-dac-trung-cua-mau-so-lieu-khong-ghep-nhom-toan-10-chuong-trinh-moi.pdf (GV. Phan Nhật Linh) (P02), Bài 02 Dạng 1 Bài tập tự luận 8"
%% kien_thuc: "Lớp 10 Bài 13 (trung vị)"
%% trang_thai: chinh-thuc
%% ghi_chu: "Tự duyệt theo yêu cầu giáo viên 22/10/2026."
\begin{ex}
Nhiệt độ ($^\circ$C) trung bình các ngày trong một tuần của một thành phố là:
\begin{center}
\begin{tabular}{|c|c|c|c|c|c|c|}\hline
$40$ & $35$ & $36$ & $39$ & $36$ & $31$ & $33$\\\hline
\end{tabular}
\end{center}
Tìm trung vị của mẫu số liệu.
\loigiai{Sắp xếp: $31$; $33$; $35$; $36$; $36$; $39$; $40$. Có $7$ giá trị nên $M_e$ là giá trị thứ $4$: $M_e=36\,^\circ$C.}
\end{ex}
